\section{Introduction}\label{sec:introduction}

For an integer \(R\ge1\), let
\[
 \Lambda_R=[-R,R]^2\cap\Z^2,\qquad
 E_R=\{\{x,y\}\subset\Lambda_R:\abs{x-y}=1\}.
\]
The planar XY model at inverse temperature \(b>0\), with free boundary,
is the probability measure
\begin{equation}\label{eq:xy-law}
 \frac{1}{Z_{R,b}}
 \exp\left\{b\sum_{\{x,y\}\in E_R}\cos(\theta_x-\theta_y)\right\}
 \prod_{x\in\Lambda_R}\frac{\dd\theta_x}{2\pi},
 \qquad \theta_x\in\R/(2\pi\Z).
\end{equation}
Each unordered edge is counted once, and there are no exterior edges.
Write \(\E^{\rm free}_{R,b}\) for expectation.  For \(r\in\Z_{>0}\), with
\(e_1=(1,0)\), set
\begin{equation}\label{eq:critical-definition}
 \begin{split}
 C_b(r)&=\lim_{\substack{R\to\infty\\R\ge r}}
       \E^{\rm free}_{R,b}\cos(\theta_0-\theta_{re_1}),\\
 m(b)&=\lim_{r\to\infty}-\frac1r\log C_b(r),\qquad
 b_c=\inf\{b>0:m(b)=0\}.
 \end{split}
\end{equation}
The limits exist, \(0<b_c<\infty\), and \(m(b_c)=0\); see
\cite[Section~1.1 and Proposition~1.2]{C}.  The thermodynamic limit in the
first line is always taken before the separation tends to infinity.

\begin{theorem}\label{thm:main}
There is a constant \(B_{\rm XY}\in(0,\infty)\) such that
\[
 C_{b_c}(r)
 =B_{\rm XY}\,r^{-1/4}(\log r)^{1/8}(1+o(1))
 \qquad(r\to\infty,\ r\in\Z).
\]
\end{theorem}

The power \(r^{-1/4}\) is the critical decay rate.  The logarithm is
smaller than every fixed power of \(r\), but it prevents
\(r^{1/4}C_{b_c}(r)\) from converging to a finite nonzero limit.  The
theorem establishes this correction and the existence of a positive
finite microscopic amplitude for the free-box thermodynamic correlation.

\subsection{Background and relation to earlier work}

Berezinskii's spin-wave analysis predicted algebraic correlations for
planar spins at low temperature \cite{Berez1971}.  Kosterlitz and
Thouless described a transition driven by the unbinding of vortex
pairs \cite{KT1973}, and Kosterlitz developed a renormalization-group
analysis of the critical behavior of the nearest-neighbor model
\cite{Kosterlitz1974}.  For the Villain formulation, Janke records the
critical prediction, up to a multiplicative constant, as
\[
 |x|^{-1/4}(\log|x|)^{1/8}
 \left(1+O\left(\frac{\log\log|x|}{\log|x|}\right)\right);
\]
see \cite[Section~II, Equations~(2)--(3)]{Janke1997}.
Theorem~\ref{thm:main} establishes the leading equivalent, with a relative
\(o(1)\) error, for the ordinary cosine interaction and the free-box
thermodynamic correlation in Equation~\eqref{eq:critical-definition}.

Rigorous work established the phase picture through correlation
bounds and dual representations.  McBryan and Spencer proved
power-law upper bounds for two-dimensional planar-spin correlations
\cite{McBryanSpencer1977}.  Fr\"ohlich and Spencer proved a
low-temperature algebraic regime for the plane rotator and Villain
models through a multiscale Coulomb-gas analysis
\cite[Theorem~C]{FrohlichSpencer1981}.  For the cosine XY model,
van Engelenburg and Lis proved exponential decay below a threshold
and the lower bound \(1/(8|x-y|)\) at and above that threshold in
the limit over free boxes \cite[Theorem~1]{vanEngelenburgLis2023}.
These results include sharp information about the phases, while their
stated bounds do not determine the critical equivalent.

The dual height field provides another route to the transition.
Aizenman, Harel, Peled, and Shapiro related height depinning to
nonsummable spin correlations, including the Bessel height interaction
dual to the cosine XY model \cite{AHPS2022}.  Lammers identified an
exact relation between the XY and dual-height correlation lengths
\cite[Theorem~1]{Lammers2023}, and proved a localization dichotomy
with a universal positive lower bound on logarithmic height
fluctuations in the delocalized phase \cite{Lammers2026}.
These results constrain the transition through height information.

The lattice-Gaussian inequality of Regev and Stephens-Davidowitz
\cite{RSD2017} underlies the finite comparisons developed in
\cite[Section~2]{B} and passed to Bessel weights in \cite[Lemma~2.2]{C}.  In a related approach,
Aizenman, Harel, Peled, and Shapiro extended sublattice monotonicity of
height moment-generating functions to annealed Gaussian interactions;
their application to the XY dual uses a representation credited there
to Raoufi \cite[Lemma~9.6, Section~10.1, and Appendix~D]{AHPS2022}.
Bauerschmidt developed finite-range decompositions of Gaussian
fields using functional calculus \cite{Bauerschmidt2013}.
Bauerschmidt, Park, and Rodriguez obtained Gaussian scaling limits
for discrete Gaussian heights at high temperature
\cite{BPR2024I,BPR2024II}.  Falco's critical Coulomb-gas analysis gives
a rigorous example in which a marginal recursion produces a power law
with a multiplicative logarithm \cite{Falco2013}.  Its fractional
electric-charge observable is different from the magnetic XY insertion
considered here.

Theorem~\ref{thm:main} sharpens the companion result
\(C_{b_c}(r)=r^{-1/4+o(1)}\) of \cite{C}.  We use from that paper
the critical height coefficient \(a(b_c)=8\pi\), finite Bessel
comparisons, fixed-geometry and conditional height limits, and a
four-face characteristic estimate.  The analytic local maps come from
\cite{B}, and the sector costs for fixed positive holes and annular kernel
estimates from \cite{S}.  The annular argument is adapted to the pinned
rim below.  The cutoff initialization, quantitative coefficient estimate,
comparison of different starting histories, and extraction of the amplitude
are proved here.

\subsection{The route to the amplitude}

Fourier expansion of the spin interactions represents a correlation
as a ratio of partition functions for integer-valued heights on the
dual lattice.  The numerator shifts the height increments by a connection
whose circulation records the spin insertion.  This magnetic insertion
is not determined by convergence of height averages alone.

For an integer \(n\ge1\), let
\begin{equation}\label{eq:center-definition}
 M_n=\E^0_{n,b_c}\cos\theta_0,
\end{equation}
where \(\E^0_{n,b_c}\) denotes expectation for the measure obtained
from Equation~\eqref{eq:xy-law} on \(\Lambda_n\) by fixing all boundary
angles at zero and integrating only the interior angles.  The geometric
argument proves two comparisons:
\[
 \begin{aligned}
 \frac{M_{n'}}{M_n}&\longrightarrow d^{-1/8}
 &&\text{if }n,n'\to\infty\text{ and }n'/n\to d\in(0,\infty),\\
 \frac{C_{b_c}(r)}{M_r^2}&\longrightarrow c_{\rm shape}\in(0,\infty)
 &&\text{as }r\to\infty.
 \end{aligned}
\]
Annular interfaces separate the microscopic insertions from an exterior
with holes kept at fixed positive radius.  The sector and annular-kernel
estimates of \cite{S} control these interfaces, and the common microscopic
factors cancel in the comparisons.  A pinned-rim argument relates the
two-insertion calculation to the free-box thermodynamic correlation.
The first comparison gives \(M_n=n^{-1/8+o(1)}\), which will make the
cutoff loss negligible relative to the center observable, and later
extends an asymptotic from geometric sizes to all integers.  The second
reduces the theorem to a center asymptotic.

To compute \(M_n\), we cut off differences of noisy cell averages,
choosing the cell size according to the target square.  For the squares
used in the analytic comparison, the rough center bound makes this change
negligible in the center partition ratio.  After averaging to the associated
starting scale, the bulk interaction is small.  For the center insertion,
every expansion term carries one labelled component that retains the source
and its attached dependencies.  The norm of the resulting labelled activity
may grow polynomially with the cell size; a linear propagation estimate
controls it.

Let \(L>1\) be the fixed dyadic block ratio and \(j\) the absolute scale index.
After separating the Gaussian reference with height coefficient \(a=8\pi\),
the smoothed bulk interaction at block side \(L^j\) retains \(t_j/2\) times
the allocated gradient energy and a fundamental periodic term; the linear
part of the remaining activity's evolution contracts.  Coordinates adapted
to the pure Gaussian curve remove its drift and give a gradient coordinate
\(u_j\) and a signed fundamental Fourier coordinate \(Z_j\).  Their leading
recursion is
\[
 u_{j+1}-u_j\simeq-HZ_j^2,\qquad
 Z_{j+1}-Z_j\simeq-2(\log L)u_jZ_j,\qquad H>0.
\]
Physical torus variance and phase tests keep each finite history near the
Gaussian curve.  A phase lower bound derived from the four-face estimate
of \cite{C} prevents the fundamental direction from disappearing.
Comparing histories through the same finite physical torus tests controls
changes of starting scale without a rate for the height scaling limit.
Together with the recursion, these constraints yield the critical balance
\(HZ_j^2\sim2(\log L)u_j^2\), which gives harmonic decay of \(u_j\).
On the controlled interior ranges of the size-dependent histories used in
the comparison, the coordinate changes give \(t_j=u_j+O(j^{-2})\).  The
resulting estimate, uniform on each fixed such range, is
\[
 t_j=\frac{1}{2(\log L)j}
       +O\left(\frac{\log j}{j^2}\right).
\]

Let \(g_n\) be the minimum squared edge norm of a connection with the
\(2\pi\) circulation of one insertion in the corresponding dual square.
For \(n_0=L^N\) and \(n_1=L^{N+1}\), we use the same entry and a common
stop \(j'=N-O(\log N)\), leaving a number of terminal blocks that grows
as a small power of \(N\).  We choose the field-independent factors removed
at each step so that their complete contributions to the insertion ratio
agree between sizes.  The source term remaining at the stop is controlled
relative to the source coefficient already extracted.  Integration of the
common phase then removes the linear periodic term, and a comparison of the
terminal Gaussian fields leaves the gradient contribution with a summable
error:
\[
 \log\frac{M_{n_1}}{M_{n_0}}
 =-\frac{1-t_{j'}}{2a}\bigl(g_{n_1}-g_{n_0}\bigr)+\epsilon_N,
 \qquad \sum_N|\epsilon_N|<\infty.
\]
Here \(g_{L^N}=2\pi(\log L)N+c_N\) with \(c_N\) convergent.  Substitution
of the trajectory estimate gives the constant increment \(-(\log L)/8\)
and the correction \(1/(16N)\).  The unweighted Green differences
\(c_{N+1}-c_N\) telescope, and summation by parts treats the differences
weighted by \(1/N\); the other errors are absolutely summable.  Summation
therefore gives a positive finite center amplitude with logarithmic power
\(1/16\).  The first geometric comparison extends this equivalent to all
integer sizes, and the bulk comparison squares the logarithmic power to
\(1/8\).
